\section{2026 Algebra \& Number Theory}

\subsection*{Exam}

\begin{exercise}
\label{algebra:2026:1}
Let $p$ be a prime, $n \geq 1$ an integer with $\gcd(n,p)=1$, and write
\[
\Gamma_0(p)=
\left\{
\begin{pmatrix}
a & b\\
c & d
\end{pmatrix}
\in \mathrm{SL}_2(\mathbb{Z})\ \middle|\ c\equiv 0 \pmod p
\right\},
\qquad
\alpha=
\begin{pmatrix}
1 & 0\\
0 & n
\end{pmatrix}.
\]
Prove that the double coset $\Gamma_0(p)\alpha\Gamma_0(p)$ decomposes into a disjoint union of right cosets
\[
\Gamma_0(p)\alpha\Gamma_0(p)=\coprod_{\gamma\in R}\Gamma_0(p)\gamma,
\]
where
\[
R=
\left\{
\begin{pmatrix}
a & b\\
0 & d
\end{pmatrix}
\ \middle|
ad=n,\ a>0,\ 0\leq b<d,\ \gcd(a,b,d)=1
\right\}.
\]
\end{exercise}

\begin{solution}
Let
\[
\Delta=
\left\{
M=
\begin{pmatrix}
A & B\\
C & D
\end{pmatrix}
\in M_2(\mathbb{Z})
\ \middle|
\det M=n,\ \gcd(A,B,C,D)=1,\ C\equiv 0 \pmod p
\right\}.
\]
We first record two elementary facts about $\Delta$.

First, $\Gamma_0(p)\alpha\Gamma_0(p)\subseteq \Delta$. Indeed, multiplication on the left and on the right by elements of $\mathrm{SL}_2(\mathbb{Z})$ preserves the determinant and the ideal generated by the four entries. Also, if
\[
g=
\begin{pmatrix}
a & b\\
c & d
\end{pmatrix},
\qquad
h=
\begin{pmatrix}
r & s\\
t & u
\end{pmatrix}
\]
belong to $\Gamma_0(p)$, then the lower-left entry of $g\alpha h$ is
\[
cr+dnt,
\]
which is divisible by $p$, since $c\equiv t\equiv 0\pmod p$.

Second, every element of $\Delta$ is left $\Gamma_0(p)$-equivalent to exactly one matrix in $R$. Let
\[
M=
\begin{pmatrix}
A & B\\
C & D
\end{pmatrix}
\in \Delta
\]
and put $g_0=\gcd(A,C)>0$. Since $C\equiv 0\pmod p$ and $p\nmid n=\det M$, the integer $A$ is not divisible by $p$; hence $p\nmid g_0$ and $C/g_0$ is divisible by $p$. Choose $r,s\in \mathbb{Z}$ with
\[
r(A/g_0)+s(C/g_0)=1.
\]
Then
\[
\eta=
\begin{pmatrix}
r & s\\
-C/g_0 & A/g_0
\end{pmatrix}
\in \Gamma_0(p),
\]
and $\eta M$ has lower-left entry $0$. Thus
\[
\eta M=
\begin{pmatrix}
g_0 & B'\\
0 & n/g_0
\end{pmatrix}
\]
for some $B'\in \mathbb{Z}$. Multiplying on the left by
\[
\begin{pmatrix}
1 & q\\
0 & 1
\end{pmatrix}
\in \Gamma_0(p)
\]
replaces $B'$ by $B'+q(n/g_0)$, so there is a unique residue $b$ with $0\leq b<n/g_0$. The resulting matrix lies in $R$, because the gcd of its entries is still $1$. Uniqueness is the usual uniqueness of Hermite normal form: if two upper triangular matrices
\[
\begin{pmatrix}
a & b\\
0 & d
\end{pmatrix},
\qquad
\begin{pmatrix}
a' & b'\\
0 & d'
\end{pmatrix}
\]
with positive diagonal entries and $0\leq b<d$, $0\leq b'<d'$ are left-equivalent under $\mathrm{SL}_2(\mathbb{Z})$, then they are the same Hermite normal form of the same row lattice. Hence $a=a'$, $d=d'$, and $b=b'$.

It remains to show that every $\gamma\in R$ actually belongs to the double coset $\Gamma_0(p)\alpha\Gamma_0(p)$. Write
\[
\gamma=
\begin{pmatrix}
a & b\\
0 & d
\end{pmatrix},
\qquad
ad=n,
\qquad
\gcd(a,b,d)=1.
\]
Choose an integer $x$ such that
\[
\gcd(ax+bp,dp)=1.
\]
This is possible by the Chinese remainder theorem: for each prime $\ell\mid d$, if $\ell\nmid a$ avoid the single residue $x\equiv -a^{-1}bp\pmod \ell$, while if $\ell\mid a$ then $\ell\nmid b$ and there is no restriction; also choose $x\not\equiv 0\pmod p$. Pick $y,w\in \mathbb{Z}$ with $xw-py=1$, and set
\[
v=
\begin{pmatrix}
x & y\\
p & w
\end{pmatrix}
\in \Gamma_0(p).
\]
The first column of $\gamma v$ is
\[
\begin{pmatrix}
ax+bp\\
dp
\end{pmatrix},
\]
which is primitive and has second entry divisible by $p$. Choose $r,s\in \mathbb{Z}$ with
\[
r(ax+bp)+sdp=1
\]
and put
\[
u=
\begin{pmatrix}
r & s\\
-dp & ax+bp
\end{pmatrix}
\in \Gamma_0(p).
\]
Then
\[
u\gamma v=
\begin{pmatrix}
1 & t\\
0 & n
\end{pmatrix}
\]
for some $t\in \mathbb{Z}$. Since
\[
\begin{pmatrix}
1 & t\\
0 & n
\end{pmatrix}
\begin{pmatrix}
1 & -t\\
0 & 1
\end{pmatrix}
=
\begin{pmatrix}
1 & 0\\
0 & n
\end{pmatrix}
=\alpha
\]
and the second factor is in $\Gamma_0(p)$, we get $\gamma\in \Gamma_0(p)\alpha\Gamma_0(p)$.

Combining the three parts, the matrices in $R$ are precisely the distinct right-coset representatives occurring in the double coset, so the asserted disjoint decomposition follows.
\end{solution}

\begin{exercise}
\label{algebra:2026:2}
Let $p$ be a prime, and let $\mathbb{F}_p$ denote the field of $p$ elements. Consider the equation
\[
Y^2=X^3+1
\]
over $\mathbb{F}_p$, and denote by $N_p$ its number of solutions in $\mathbb{F}_p^2$:
\[
N_p:=\#\{(x,y)\in \mathbb{F}_p^2\mid y^2=x^3+1\}.
\]
In the following, we want to show the inequality
\[
|N_p-p|\leq 2\sqrt{p}.
\]
For $m\in \mathbb{Z}_{\geq 1}$, let $X_m=X_m(\mathbb{F}_p^\times)$ be the group of complex characters of $\mathbb{F}_p^\times$ of order dividing $m$. For a character $\chi$ of $\mathbb{F}_p^\times$, set
\[
\chi(0)=
\begin{cases}
1, & \text{if }\chi=1,\\
0, & \text{otherwise.}
\end{cases}
\]
For two characters $\chi$ and $\mu$ of $\mathbb{F}_p^\times$, write $J(\chi,\mu)$ for the Jacobi sum attached to them:
\[
J(\chi,\mu):=\sum_{a+b=1}\chi(a)\mu(b)\in \mathbb{C}.
\]
Recall that $|J(\chi,\mu)|=\sqrt{p}$ if the characters $\chi$, $\mu$, and $\chi\mu$ are all non-trivial.

\begin{enumerate}
\item[(1)] Assume $p=3$ or $p\equiv 2\pmod 3$. Show that $N_p=p$, and thus $|N_p-p|\leq 2\sqrt{p}$ in this case.
\item[(2)] Assume $p\equiv 1\pmod m$. Let $a\in \mathbb{F}_p$. Write $N(X^m=a)$ for the number of solutions of the equation $X^m=a$. Show that
\[
N(X^m=a)=\sum_{\chi\in X_m}\chi(a).
\]
\item[(3)] Suppose $p\equiv 1\pmod 3$. Show that we still have $|N_p-p|\leq 2\sqrt{p}$.
\end{enumerate}
\end{exercise}

\begin{solution}
\textbf{Part (1).}
If $p=3$, then $x^3=x$ for every $x\in \mathbb{F}_3$. If $p\equiv 2\pmod 3$, then $\gcd(3,p-1)=1$, so the map $x\mapsto x^3$ is a bijection of $\mathbb{F}_p^\times$ and hence of $\mathbb{F}_p$. Therefore, in both cases $x\mapsto x^3$ is a bijection of $\mathbb{F}_p$.

Thus the number of pairs $(x,y)$ satisfying $y^2=x^3+1$ is the same as the number of pairs $(t,y)$ satisfying $y^2=t+1$. For each $y\in \mathbb{F}_p$ there is exactly one such $t$, namely $t=y^2-1$. Hence $N_p=p$.

\textbf{Part (2).}
Assume $p\equiv 1\pmod m$. Then $\mathbb{F}_p^\times$ is cyclic of order divisible by $m$, and $X_m$ has exactly $m$ characters.

If $a=0$, then $X^m=0$ has the unique solution $X=0$. On the right-hand side only the trivial character contributes at $0$, so
\[
\sum_{\chi\in X_m}\chi(0)=1.
\]
Thus the formula holds for $a=0$.

Now let $a\in \mathbb{F}_p^\times$. The character orthogonality relation on the finite cyclic group $\mathbb{F}_p^\times/(\mathbb{F}_p^\times)^m$ gives
\[
\sum_{\chi\in X_m}\chi(a)
=
\begin{cases}
m, & a\in (\mathbb{F}_p^\times)^m,\\
0, & a\notin (\mathbb{F}_p^\times)^m.
\end{cases}
\]
The homomorphism $x\mapsto x^m$ on $\mathbb{F}_p^\times$ has kernel of order $m$, so $X^m=a$ has exactly $m$ solutions if $a$ is an $m$-th power and no solution otherwise. This proves the formula.

\textbf{Part (3).}
Assume $p\equiv 1\pmod 3$. Then $p$ is odd. Let $\varphi$ be the quadratic character of $\mathbb{F}_p^\times$, extended to $0$ by $\varphi(0)=0$, and let $\rho$ be a non-trivial cubic character. Then
\[
X_3=\{1,\rho,\rho^2\}.
\]
For each $A\in \mathbb{F}_p$, the equation $y^2=A$ has $1+\varphi(A)$ solutions. Using Part (2) for cubes, we get
\[
N_p
=
\sum_{t\in \mathbb{F}_p}
\left(1+\rho(t)+\rho^2(t)\right)
\left(1+\varphi(t+1)\right).
\]
The sums of the non-trivial characters $\rho$, $\rho^2$, and $\varphi(t+1)$ over $\mathbb{F}_p$ are all $0$. Therefore
\[
N_p-p
=
\sum_{t\in \mathbb{F}_p}\rho(t)\varphi(t+1)
+
\sum_{t\in \mathbb{F}_p}\rho^2(t)\varphi(t+1).
\]
Make the change of variables $u=-t$. Since $p\equiv 1\pmod 3$, the element $-1$ is a cube in $\mathbb{F}_p^\times$, and hence $\rho(-1)=\rho^2(-1)=1$. Therefore
\[
\sum_{t\in \mathbb{F}_p}\rho(t)\varphi(t+1)
=
\sum_{u\in \mathbb{F}_p}\rho(u)\varphi(1-u)
=
J(\rho,\varphi),
\]
and
\[
\sum_{t\in \mathbb{F}_p}\rho^2(t)\varphi(t+1)
=
\sum_{u\in \mathbb{F}_p}\rho^2(u)\varphi(1-u)
=
J(\rho^2,\varphi).
\]
The characters $\rho$, $\varphi$, $\rho\varphi$ are all non-trivial, and the characters $\rho^2$, $\varphi$, $\rho^2\varphi$ are also all non-trivial. Hence the recalled Jacobi-sum estimate gives
\[
|J(\rho,\varphi)|=\sqrt p,
\qquad
|J(\rho^2,\varphi)|=\sqrt p.
\]
Consequently
\[
|N_p-p|
\leq
|J(\rho,\varphi)|+|J(\rho^2,\varphi)|
=2\sqrt p.
\]
\end{solution}

\begin{exercise}
\label{algebra:2026:3}
Let $K$ be a number field, with $\mathcal{O}_K$ its ring of integers. Let $\overline K$ be an algebraic closure of $K$, and denote by $\mathcal{O}_{\overline K}$ the set of algebraic integers, i.e.\ the set of elements of $\overline K$ that are integral over $\mathcal{O}_K$. Let $a,b\in \mathcal{O}_K$. Show that the following two conditions are equivalent:
\begin{enumerate}
\item[(1)] the two elements $a$ and $b$ are coprime to each other, in the sense that the ideal $(a,b)\subset \mathcal{O}_K$ generated by $a,b$ is $\mathcal{O}_K$;
\item[(2)] there exists some $u\in \mathcal{O}_{\overline K}$ so that $au+b\in \mathcal{O}_{\overline K}^{\times}$.
\end{enumerate}
\end{exercise}

\begin{solution}
Put $A=\mathcal{O}_K$ and $B=\mathcal{O}_{\overline K}$.

First assume that $(a,b)\ne A$. Then there is a non-zero prime ideal $\mathfrak p\subset A$ containing both $a$ and $b$. Since $B$ is integral over $A$, the lying-over theorem gives a prime ideal $\mathfrak P\subset B$ with $\mathfrak P\cap A=\mathfrak p$. For any $u\in B$ we have
\[
au+b\in \mathfrak P,
\]
because $a,b\in \mathfrak P$. A proper prime ideal contains no unit, so $au+b$ cannot lie in $B^\times$. Thus (2) implies (1).

Conversely assume $(a,b)=A$. If $a=0$, then $b\in A^\times$, and $u=0$ works. If $a\in A^\times$, choose
\[
u=a^{-1}(1-b)\in A;
\]
then $au+b=1$. Hence we may assume that $a$ is non-zero and not a unit.

Because $(a,b)=A$, the image of $b$ is a unit in the finite ring $A/(a)$. Hence the image of $-b$ has finite multiplicative order. Choose $N\geq 1$ such that
\[
(-b)^N\equiv 1\pmod a.
\]
Then $1-(-b)^N\in (a)$. The ideal generated by the elements
\[
a(-b)^{N-1},\ a^2(-b)^{N-2},\ \ldots,\ a^N
\]
is
\[
a(a,b)^{N-1}=(a).
\]
Therefore there exist $c_{N-1},c_{N-2},\ldots,c_0\in A$ such that
\[
1-(-b)^N
=
\sum_{i=1}^N c_{N-i}(-b)^{N-i}a^i.
\]
Define the monic polynomial
\[
F(T)=T^N+c_{N-1}T^{N-1}+\cdots+c_1T+c_0\in A[T].
\]
The preceding identity is exactly
\[
a^N F(-b/a)=1.
\]
Let $u\in \overline K$ be any root of $F$. Since $F$ is monic over $A$, the element $u$ belongs to $B$.

Let $u_1,\ldots,u_N$ be the roots of $F$ in $\overline K$, counted with multiplicity. Then
\[
a^N F(-b/a)
=
(-1)^N\prod_{i=1}^N (b+au_i).
\]
The left-hand side is $1$, so the product $\prod_i(b+au_i)$ is $\pm 1$. Each factor $b+au_i$ is an algebraic integer. A factor of a unit among algebraic integers is again a unit: explicitly, the inverse of $b+au_i$ is $\pm$ the product of all the remaining factors. Hence $b+au_i\in B^\times$ for every $i$. Taking $u=u_1$ proves (2).
\end{solution}

\begin{exercise}
\label{algebra:2026:4}
Let $a\in \mathbb{Z}$ be square-free and let $\alpha=\sqrt[3]{a}$ be a root of $f(x)=x^3-a$. Recall that the discriminant $\mathrm{Disc}(f)=-27a^2$. Denote $K=\mathbb{Q}(\alpha)$. Show that the ring of integers $\mathcal{O}_K$ has integral basis
\[
\{1,\alpha,\alpha^2\}
\quad\text{if }a\not\equiv \pm 1\pmod 9,
\]
and
\[
\left\{1,\alpha,\frac{1\pm \alpha+\alpha^2}{3}\right\}
\quad\text{if }a\equiv \pm 1\pmod 9.
\]
\end{exercise}

\begin{solution}
We prove the usual non-degenerate pure cubic case, namely $a\ne \pm 1$, so that $x^3-a$ is irreducible and $[K:\mathbb{Q}]=3$. The cases $a=\pm 1$ are degenerate: then $K=\mathbb{Q}$, so the displayed three-element bases are not bases of $K$.

The power basis $\{1,\alpha,\alpha^2\}$ has discriminant
\[
\mathrm{Disc}(1,\alpha,\alpha^2)=\mathrm{Disc}(x^3-a)=-27a^2.
\]
Let
\[
I=[\mathcal{O}_K:\mathbb{Z}[\alpha]].
\]
Then $I^2$ divides $27a^2$. If a prime $\ell\neq 3$ divides $a$, then $x^3-a$ is Eisenstein at $\ell$, since $a$ is square-free; hence $\ell\nmid I$. If $\ell\neq 3$ does not divide $a$, then $\ell\nmid \mathrm{Disc}(x^3-a)$, so $\ell\nmid I$. If $3\mid a$, then $x^3-a$ is Eisenstein at $3$, again because $a$ is square-free, and $3\nmid I$. Thus the only remaining case is $3\nmid a$, and the only possible prime divisor of $I$ is $3$.

Assume $3\nmid a$. Let $\varepsilon\in\{1,-1\}$ be the sign with
\[
a\equiv \varepsilon\pmod 3.
\]
Modulo $3$ we have
\[
x^3-a\equiv (x-\varepsilon)^3\pmod 3.
\]
Use Dedekind's index criterion with the lift $\phi(x)=x-\varepsilon$. Since
\[
(x-\varepsilon)^3=x^3-3\varepsilon x^2+3x-\varepsilon,
\]
we get
\[
\frac{x^3-a-(x-\varepsilon)^3}{3}
=
\varepsilon x^2-x+\frac{\varepsilon-a}{3}.
\]
Evaluating this polynomial at $x=\varepsilon$ gives
\[
\frac{\varepsilon-a}{3}.
\]
Therefore $3$ divides the index $I$ exactly when
\[
\frac{\varepsilon-a}{3}\equiv 0\pmod 3,
\]
that is, exactly when $a\equiv \varepsilon\pmod 9$, equivalently $a\equiv \pm 1\pmod 9$. Since $v_3(-27a^2)=3$ in the case $3\nmid a$, the index can then be only $1$ or $3$. Hence
\[
I=
\begin{cases}
1, & a\not\equiv \pm 1\pmod 9,\\
3, & a\equiv \pm 1\pmod 9.
\end{cases}
\]
This proves the first asserted basis.

It remains to identify the overorder in the second case. Suppose $a\equiv \varepsilon\pmod 9$, where $\varepsilon=\pm 1$, and put
\[
\beta=\frac{1+\varepsilon\alpha+\alpha^2}{3}.
\]
We show that $\beta$ is integral. From
\[
3\beta=1+\varepsilon\alpha+\alpha^2
\]
one obtains, by eliminating $\alpha$ using $\alpha^3=a$, the equation
\[
\beta^3-\beta^2-\frac{\varepsilon(a-\varepsilon)}{3}\beta-\frac{(a-\varepsilon)^2}{27}=0.
\]
The coefficients are integers because $a\equiv \varepsilon\pmod 9$. Thus $\beta$ is an algebraic integer.

The change-of-basis determinant from $\{1,\alpha,\alpha^2\}$ to $\{1,\alpha,\beta\}$ is $1/3$. Therefore
\[
\mathrm{Disc}(1,\alpha,\beta)
=
\frac{\mathrm{Disc}(1,\alpha,\alpha^2)}{9}
=
-3a^2.
\]
Since the index of $\mathbb{Z}[\alpha]$ in $\mathcal{O}_K$ is exactly $3$, adjoining this integral element accounts for the whole index. Hence
\[
\mathcal{O}_K=\mathbb{Z}\oplus \mathbb{Z}\alpha\oplus \mathbb{Z}\beta.
\]
Equivalently, if $a\equiv 1\pmod 9$ then the basis is
\[
\left\{1,\alpha,\frac{1+\alpha+\alpha^2}{3}\right\},
\]
and if $a\equiv -1\pmod 9$ then the basis is
\[
\left\{1,\alpha,\frac{1-\alpha+\alpha^2}{3}\right\}.
\]
\end{solution}

\begin{exercise}
\label{algebra:2026:5}
Let $p$ be a prime number. Let $K/\mathbb{Q}_p$ be a finite extension of $\mathbb{Q}_p$, the field of $p$-adic numbers.
\begin{enumerate}
\item[(1)] Show that for any fixed integer $d\geq 1$, $K$ has only finitely many non-isomorphic finite extensions of degree $d$.
\item[(2)] Assume $(d,p)=1$. Suppose that $K$ has a finite Galois extension $L/K$ of degree $d$ which is totally ramified. Show that $K$ contains a primitive $d$-th root of unity and that $L/K$ must be cyclic.
\item[(3)] Assume $(d,p)=1$. Compute the number of non-isomorphic Galois totally ramified extensions of $K$ of degree $d$.
\end{enumerate}
\end{exercise}

\begin{solution}
Let $\mathcal{O}_K$ be the ring of integers of $K$, let $\mathfrak p_K$ be its maximal ideal, let $k$ be its residue field, and put $q=\#k$.

\textbf{Part (1).}
We use the standard compactness proof based on Krasner's lemma. First consider totally ramified extensions of a fixed degree $e$. Every such extension is generated by a uniformizer, whose minimal polynomial over $K$ is Eisenstein of degree $e$. The set of Eisenstein polynomials
\[
T^e+c_{e-1}T^{e-1}+\cdots+c_1T+c_0
\]
with $c_i\in \mathfrak p_K$ for $i\geq 1$ and $c_0\in \mathfrak p_K\setminus \mathfrak p_K^2$ is a compact subset of $\mathcal{O}_K^e$. By Krasner's lemma, sufficiently small perturbations of a fixed Eisenstein polynomial generate the same degree-$e$ extension. Thus this compact set is covered by open subsets on each of which the generated extension is constant; a finite subcover gives only finitely many totally ramified extensions of degree $e$.

For a general extension $E/K$ of degree $d$, let $E_0/K$ be its maximal unramified subextension. Then $[E_0:K]=f$ divides $d$, and for each $f$ there is a unique unramified extension of $K$ of degree $f$. The extension $E/E_0$ is totally ramified of degree $d/f$. Applying the totally ramified finiteness just proved over each of the finitely many possible fields $E_0$, we obtain only finitely many degree-$d$ extensions of $K$.

\textbf{Part (2).}
Assume $(d,p)=1$ and let $L/K$ be Galois, totally ramified, and of degree $d$. Since the ramification index is prime to $p$, the extension is tame. Choose a uniformizer $\pi_L$ of $L$. For $\sigma\in \mathrm{Gal}(L/K)$, the quotient
\[
\frac{\sigma(\pi_L)}{\pi_L}
\]
is a unit of $\mathcal{O}_L$. Reducing modulo the maximal ideal gives an element of $k_L^\times=k^\times$, because $L/K$ is totally ramified.

This defines the tame inertia character
\[
\theta:\mathrm{Gal}(L/K)\longrightarrow k^\times,
\qquad
\theta(\sigma)=\overline{\sigma(\pi_L)/\pi_L}.
\]
It is injective in the tame totally ramified case. Indeed, if $\theta(\sigma)=1$, then $\sigma$ acts trivially on both the residue field and on $\pi_L/\pi_L^2$; the usual tame ramification filtration argument gives $\sigma=1$. Hence $\mathrm{Gal}(L/K)$ embeds into the cyclic group $k^\times$. Therefore $\mathrm{Gal}(L/K)$ is cyclic and $d\mid q-1$.

Since $d\mid q-1$, the finite field $k$ contains a primitive $d$-th root of unity. Its Teichmuller lift is a primitive $d$-th root of unity in $K$. Thus $K$ contains $\mu_d$, and $L/K$ is cyclic.

\textbf{Part (3).}
If $d\nmid q-1$, then by Part (2) no such extension exists. Assume now that $d\mid q-1$, equivalently $\mu_d\subset K$.

We use Kummer theory only in its tame local form: because $(d,p)=1$ and $\mu_d\subset K$, cyclic extensions of exponent dividing $d$ correspond to cyclic subgroups of $K^\times/K^{\times d}$ via
\[
c\longmapsto K(\sqrt[d]{c}).
\]

Because $(d,p)=1$, the map $x\mapsto x^d$ is an automorphism of the principal unit group $1+\mathfrak p_K$ by Hensel's lemma applied to $T^d-u$. Choosing a uniformizer $\pi$ of $K$ and a Teichmuller lift $\omega\in \mathcal{O}_K^\times$ whose residue class generates $k^\times/(k^\times)^d$, we have
\[
K^\times/K^{\times d}
\cong
\langle \pi\rangle/\langle \pi^d\rangle
\oplus
k^\times/k^{\times d}
\cong
(\mathbb{Z}/d\mathbb{Z})^2.
\]
Under this identification, write the class of $\pi^r\omega^s$ as $(r,s)$.

Such a Kummer extension is totally ramified exactly when the valuation component of the corresponding class has order $d$: the residue-field component gives the unramified part, while the valuation component gives the ramification index. Thus the cyclic subgroups of order $d$ that produce totally ramified extensions are precisely those whose first projection onto $\mathbb{Z}/d\mathbb{Z}$ is an isomorphism. These subgroups are precisely
\[
\langle (1,j)\rangle,
\qquad
j\in \mathbb{Z}/d\mathbb{Z}.
\]
Thus there are exactly $d$ non-isomorphic Galois totally ramified extensions of degree $d$ when $d\mid q-1$.

Equivalently, representatives are
\[
K\left(\sqrt[d]{\pi\omega^j}\right),
\qquad
j=0,1,\ldots,d-1.
\]
Therefore the required number is
\[
\begin{cases}
d, & d\mid q-1,\\
0, & d\nmid q-1.
\end{cases}
\]
\end{solution}
